\documentclass{resume}
% \usepackage{zh_CN-Adobefonts_external} % Simplified Chinese Support using external fonts (./fonts/zh_CN-Adobe/)
% \usepackage{zh_CN-Adobefonts_internal} % Simplified Chinese Support using system fonts

\newcommand{\SoftwareEngineer}{Software Engineer}
\newcommand{\DataScientist}{Data Scientist}
\newcommand{\ResearchFellow}{Postdoctoral Research Fellow}

% \usepackage{fontawesome5}


\usepackage{xcolor}
\begin{document}
\pagenumbering{gobble}

{\LARGE \textbf{Guanli Liu}}\\[-0.2em]
% \faEnvelope\ \href{mailto:guanlil1@unimelb.edu.au}{guanlil1@unimelb.edu.au} 
\faEnvelope\ \href{mailto:liuguanli22@gmail.com}{liuguanli22@gmail.com} 
\textperiodcentered\ 
\faLinkedin\ \href{https://www.linkedin.com/in/guanli-liu-11353058/}{LinkedIn} \textperiodcentered\ 
\faGithub\ \href{https://github.com/Liuguanli}{GitHub}

\faMapMarker\ Melbourne, Australia \textperiodcentered\ Open to relocation and remote \textperiodcentered\ 
\faClockO\  UTC+10/+11 \textperiodcentered\ 
\href{https://www.worldtimebuddy.com/?lid=5808079,5128581,5809844,1816670,2158177&h=2158177&hf=1}{Global time overlap \faExternalLink}
\vspace{0.4em}

% Redmond 是 5808079，San Francisco 是 5391959，New York City 是 5128581，Seattle 是 5809844，Singapore 是 1880252，Munich 是 2867714，Barcelona 是 3128760，Vancouver 是 6173331，Tokyo 是 1850147，Melbourne 是 2158177，而被去掉的 Beijing 是 1816670


% -------------------------
% Profile
% -------------------------
\section{\faUser\ \textbf{Profile}}
Software engineer and postdoctoral researcher with a strong background in database systems, query processing, indexing, physical design, and performance optimization. My work has focused on system internals that affect query behavior, data layout, benchmarking, and performance trade-offs. I have built systems for workload parsing, layout recommendation, and systems evaluation, and published related work in venues such as VLDB and ICDE. Across research and industry, I have also worked on backend services, database performance tuning, and production-oriented data workflows.



% -------------------------
% Research & Work Experience
% -------------------------
\section{\faUsers\ \textbf{Research and Work Experience}}

\datedsubsection{\textbf{PhD + Postdoc}, \textbf{The University of Melbourne}, Australia}{June. 2019 -- Present}

\textit{Lead research and engineering projects on database systems, physical design, indexing, and performance evaluation.}
\textit{Build system prototypes, study trade-offs in query behavior, and collaborate across research and engineering settings.}

\begin{itemize}[parsep=0.4ex]


\item \textbf{Data Layout and Physical Design}
\begin{itemize}[parsep=0.3ex]
    \item Built a \textbf{layout advisory} system for data lake-style datasets, supporting dataset and workload ingestion, SQL workload parsing, and automatic physical data optimization through layout recommendation.
    \href{https://github.com/Liuguanli/layout_advisor}{\faGithub}
    \textit{[VLDB 2026, submitted]}

    \item Proposed a file \textbf{layout method} for data lake systems and developed an \(\mathcal{O}(1)\) cost model to efficiently measure layout quality.
    \href{https://github.com/Liuguanli/LBMC}{\faGithub}
    \href{https://www.vldb.org/pvldb/vol17/p4773-liu.pdf}{\textit{[VLDB 2025\faExternalLink]}}
\end{itemize}

\item \textbf{Benchmarking and Evaluation}
\begin{itemize}[parsep=0.3ex]
    \item Built \textbf{DriftBench}, a framework for evaluating performance stability under workload and data drift across database tuning and benchmarking scenarios.
    \href{https://github.com/Liuguanli/DriftBench}{\faGithub}
    \href{https://arxiv.org/abs/2510.10858}{\textit{[VLDB 2026\faExternalLink]}}
    \item Built a \textbf{benchmarking framework} for traditional and learned spatial indexes, analyzing latency, I/O, and workload sensitivity under controlled query workloads. \href{https://github.com/Liuguanli/rl_spatial_benchmark}{\faGithub}
\href{https://arxiv.org/abs/2512.11161}{\textit{[ICDE 2026\faExternalLink]}}
\end{itemize}

\item \textbf{Indexing, Spatial Indexes and Learned Indexes}
\begin{itemize}[parsep=0.3ex]
    \item Proposed a learned index construction method based on model reuse and fine-tuning, improving the efficiency and adaptability of index learning.
\href{https://people.eng.unimelb.edu.au/jianzhongq/papers/DBML2023_ModelReuse.pdf}{\textit{[ICDEW 2023\faExternalLink]}}

\item Built learned spatial index structures for multi-dimensional data, enabling efficient spatial query processing with improved lookup performance.
\href{https://people.eng.unimelb.edu.au/baileyj/papers/ICDE2023_ELSI.pdf}{\textit{[ICDE 2023\faExternalLink]}}

\item Developed recursive learned spatial index structures and studied their effectiveness for spatial query workloads.
\href{https://github.com/Liuguanli/RSMI}{\faGithub}
\href{https://www.vldb.org/pvldb/vol13/p2341-qi.pdf}{\textit{[VLDB 2020\faExternalLink]}}
\end{itemize}

\item \textbf{Other (Supervised Works)}
\begin{itemize}[parsep=0.3ex]
    \item Supervised work on learned cardinality estimation models. 
    \href{<paper-link>}{\textit{[ICDE 2026: CoLSE\faExternalLink]}}

   

    \item Supervised work on LLM-assisted processing of complex spatial queries. 
    \href{https://github.com/AI-DB-UoM/complex_spatial_data_generation}{\faGithub}
    \href{https://findanexpert.unimelb.edu.au/scholarlywork/2270718-llm-enhanced-processing-of-complex-spatial-queries?cache=1772143387383}{\textit{[ADC 2025\faExternalLink]}}

     \item Supervised work on embedding-based spatial keyword query processing.{\textit{[ADC 2025]}}
\end{itemize}

\end{itemize}

\datedsubsection{\textbf{Data Scientist}, \textbf{{nftDb}}, Australia}{Feb. 2023 -- Feb. 2024}
\textit{A startup focused on NFT trading data, market trend tracking, and multimodal search over digital assets.}

\begin{itemize}[parsep=0.4ex]

\item Built Python-based ingestion pipelines for raw blockchain transaction data, supporting Kafka-based streaming and Airflow-orchestrated batch workflows for downstream analytics.

\item Developed SQL and dbt workflows to clean, normalize, and model large-scale transaction data into analytics-ready tables for product insights and reporting.

\item Queried and analyzed transaction-level datasets in BigQuery to uncover user behaviour, marketplace activity, and performance trends across digital asset platforms.

\item Developed an internal RAG-based knowledge assistant for technical documentation retrieval, data asset discovery, and workflow support across engineering teams.

\item Partnered with engineers to improve pipeline reliability, data quality, and metric consistency across the end-to-end analytics and reporting stack.

\end{itemize}

% \datedsubsection{\textbf{Research Assistant}, \textbf{The University of Melbourne}, Australia}{Aug. 2022 -- Aug. 2023}
% \begin{itemize}
%   \item \textit{Project One:} Developed an AI-assisted system for reducing reading interruptions using \textbf{GPT API} and \textbf{Google Cloud}.
%   \item \textit{Project Two:} Designed a \textbf{C language coding style checker} to detect common errors using Python.
% \end{itemize}

\datedsubsection{\textbf{Software Engineer}, \textbf{Baidu}, China}{Jul. 2015 -- Aug. 2017}

\textit{Worked on Baidu's internal company IM system serving employees and business partners.}
\begin{itemize}[parsep=0.4ex]

\item Designed messaging protocols and implemented message deduplication mechanisms to improve delivery reliability and consistency across large-scale communication workflows.

% \item Collaborated with backend engineers and cross-platform teams to design file transfer architecture supporting multi-end business scenarios and efficient data synchronization.

\item Improved database performance through profiling, query optimization, and systematic tuning of backend data access workflows.

\end{itemize}


% -------------------------
% Skills
% -------------------------
\section{\faCogs\ \textbf{Engineering Skills}}

\begin{itemize}[parsep=0.4ex]

\item \textbf{Programming:} Proficient in Python and Java, with working knowledge of C++ and SQL

\item \textbf{Database / Systems:} Query processing, indexing, physical design, benchmarking, and performance evaluation

\item \textbf{Data Systems \& Pipelines:} Experience with SparkSQL, Apache Hudi, REST APIs, data ingestion pipelines, and batch / incremental data processing

\item \textbf{Databases:} Experience with PostgreSQL, PostGIS, pgvector, BigQuery, and spatial data management

\item \textbf{AI / Retrieval Systems:} Experience with RAG pipelines, vector databases, embedding-based retrieval

\item \textbf{Cloud \& Infrastructure:} Experience with Google Cloud Platform and Docker

\item \textbf{Machine Learning:} Hands-on experience with PyTorch, TorchLib, Scikit-learn, and common machine learning algorithms for classification, regression, clustering, and representation learning

  \item \textbf{Algorithmic Knowledge:} Solid understanding of fundamental data structures and algorithms
\end{itemize}

% -------------------------
% Education
% -------------------------
\section{\faGraduationCap\ \textbf{Education}}
\datedsubsection{\textbf{PhD} in Computer Science, \textit{The University of Melbourne}, Australia}{2019 -- 2023}
\datedsubsection{\textbf{M.S.} in Computer Technology, \textit{Northeastern University}, China}{2013 --  2015}
\datedsubsection{\textbf{B.Eng.} in Software Engineering, \textit{Northeastern University}, China}{2009 --  2013}

% -------------------------
% Publications
% -------------------------
% \section{\faFile\ \textbf{Selected Publications}}

% \textit{Representative publications from projects that I led or supervised.}

% \begin{itemize}

% \item ``Toward Drift-Aware Database Benchmarking.'' \textbf{\textit{VLDB}}, 2026. \textit{[1st author]}

% \item ``Benchmarking RL-Enhanced Spatial Indices Against Traditional, Advanced, and Learned Counterparts.'' \textbf{\textit{ICDE}}, 2026. \textit{[1st author]}

% \item ``CoLSE: A Lightweight and Robust Hybrid Learned Model for Single-Table Cardinality Estimation using Joint CDF.'' \textbf{\textit{ICDE}}, 2026. \textit{[2nd author]}

% \item ``Advancing Spatial Keyword Queries: From Filters to Unified Vector Embeddings.'' \textbf{\textit{ADC}}, 2025. \textit{[2nd author]}

% \item ``LLM-Enhanced Processing of Complex Spatial Queries.'' \textbf{\textit{ADC}}, 2025. \textit{[2nd author]}

% \item ``Efficient Cost Modeling of Space-filling Curves.'' \textbf{\textit{VLDB}}, 2025. \textit{[1st author]}

% \item ``Efficient Index Learning via Model Reuse and Fine-tuning.'' \textbf{\textit{ICDEW}}, 2023. \textit{[1st author]}

% \item ``Efficiently Learning Spatial Indices.'' \textbf{\textit{ICDE}}, 2023. \textit{[1st author]}

% \item ``Effectively Learning Spatial Indices.'' \textbf{\textit{VLDB}}, 2020. \textit{[2nd author]}

% \item ``The Moving K Diversified Nearest Neighbor Query.'' \textbf{\textit{TKDE}}, 2016. \textit{[2nd author]}

% \end{itemize}
% % \section{\faFile\ \textbf{Selected Publications}}

% % \begin{itemize}

% % \item \textbf{Guanli Liu}, Renata Borovica-Gajic, Hai Lan, Zhifeng Bao.  
% % ``Benchmarking RL-Enhanced Spatial Indices Against Traditional, Advanced, and Learned Counterparts.''  
% % \textit{ICDE}, 2026. 
% % % \href{https://github.com/Liuguanli/rl_spatial_benchmark}{\textcolor{blue}{\texttt{[Code]}}}

% % \item Lankadinee Rathuwadu, \textbf{Guanli Liu}, Christopher Leckie, Renata Borovica-Gajic.  
% % ``CoLSE: A Lightweight and Robust Hybrid Learned Model for Single-Table Cardinality Estimation using Joint CDF.''  
% % \textit{ICDE}, 2026.

% % \item Kaan Gocmen, \textbf{Guanli Liu}, Renata Borovica-Gajic.  
% % ``Advancing Spatial Keyword Queries: From Filters to Unified Vector Embeddings.''  
% % \textit{Australasian Database Conference (ADC)}, 2025.

% % \item Ruiyi Hao, \textbf{Guanli Liu}, Renata Borovica-Gajic.  
% % ``LLM-Enhanced Processing of Complex Spatial Queries.''  
% % \textit{Australasian Database Conference (ADC)}, 2025.


% %   \item \textbf{Guanli Liu}, Lars Kulik, Christian S. Jensen, Tianyi Li, Renata Borovica-Gajic, Jianzhong Qi.  
% %   ``Efficient Cost Modeling of Space-filling Curves." \textit{VLDB}, 2025.  

  

% %   \item \textbf{Guanli Liu}, Jianzhong Qi, Lars Kulik, Kazuya Soga, Renata Borovica-Gajic, Benjamin I. P. Rubinstein.  
% %   ``Efficient Index Learning via Model Reuse and Fine-tuning." \textit{ICDEW}, 2023.  


% %   \item \textbf{Guanli Liu}, Jianzhong Qi, Christian S. Jensen, James Bailey, Lars Kulik.  
% %   ``Efficiently Learning Spatial Indices." \textit{ICDE}, 2023.  


% %   \item Jianzhong Qi, \textbf{Guanli Liu}, Christian S. Jensen, Lars Kulik.  
% %   ``Effectively Learning Spatial Indices." \textit{VLDB}, 2020.  

% %   \item Yu Gu, \textbf{Guanli Liu}, Jianzhong Qi, Hongfei Xu, Ge Yu, Rui Zhang.  
% %   ``The Moving K Diversified Nearest Neighbor Query." \textit{TKDE}, 2016.  

% % \end{itemize}


% % % -------------------------
% % % Seminar/Conference Presentations
% % % -------------------------
% % \section{\faMicrophone\ \textbf{Seminar and Conference Presentations}}
% % \begin{itemize}
% %   \item TODO (List conferences where you have presented, specify if invited talk, oral, or poster)
% % \end{itemize}

% % % -------------------------
% % % Teaching and Supervision
% % % -------------------------
% % \section{\faUniversity\ \textbf{Teaching and Supervision}}
% % \begin{itemize}
% %   \item TODO (List postgraduate/undergraduate teaching roles, tutor/demo roles, co-supervision experience)
% % \end{itemize}

\section{\faGraduationCap\ \textbf{Teaching}}
\begin{itemize}
  % \item \textbf{INFO20003 – Database Systems (The University of Melbourne):} 
  % Prepared teaching materials for the course in advance of the 2025 Semester 2 offering. Will participate in teaching in the upcoming semester.

  \item \textbf{COMP90018 – Android Application Development (The University of Melbourne):} 
  Tutor from Aug. 2019 to 2023, responsible for tutorials, student support, and assessment marking.

  \item \textbf{COMP90041 – Programming and Software Development (The University of Melbourne):} 
  Tutor from Aug. 2019 to 2023, responsible for tutorials and assessment marking
\end{itemize}

% -------------------------
% Leadership and Service
% -------------------------
\section{\faUsers\ \textbf{Research Service}}
\begin{itemize}
  \item \textbf{Conference Reviewer:}  \textbf{SIGMOD} 2026, \textbf{VLDB} 2027 (PC), 2026 (Shadow PC), 2025 (External), KDD 2026, 2025 (Excellent Reviewer)
  \item \textbf{Journal Reviewer:} TKDE, WWW, Transactions on Spatial Algorithms and Systems (TSAS)
\end{itemize}

% \section{\faBriefcase\ \textbf{Supervision}}


% \noindent\textbf{Kaan Gocmen} (Master’s Student, University of Melbourne)\\
% \textit{Topic:} Advancing Spatial Keyword Queries\\
% \textit{Publication:} To appear in \textit{Australasian Database Conference (ADC)}, 2025.

% \vspace{0.8em}
% \noindent\textbf{Ruiyi Hao} (Master’s Student, University of Melbourne)\\
% \textit{Topic:} LLM for Spatial Queries\\
% \textit{Publication:} To appear in \textit{Australasian Database Conference (ADC)}, 2025.

% \end{itemize}

% % -------------------------
% % Administrative Responsibilities
% % -------------------------
% \section{\faTasks\ \textbf{Administrative Responsibilities}}
% \begin{itemize}
%   \item TODO (List any faculty, department, lab, or institutional responsibilities)
% \end{itemize}

% % -------------------------
% % Community Engagement
% % -------------------------
% \section{\faUsers\ \textbf{Community Engagement}}
% \begin{itemize}
%   \item TODO (List any public talks, outreach, science communication, open-source contributions)
% \end{itemize}

\end{document}
